\noindent\textbf{Student Name:} \rule{0.45\textwidth}{0.4pt}
\hfill
\textbf{Student ID:} \rule{0.2\textwidth}{0.4pt}

\vspace{0.75em}

\subsection*{Instructions (Instructor Use)}
Circle the level that best reflects the student’s overall performance since the previous checkpoint. Use the booklet as supporting evidence, but prioritise observed participation.

\vspace{0.75em}

\renewcommand{\arraystretch}{1.4}
{\scriptsize
\begin{tabularx}{\textwidth}{|X|X|X|X|X|X|}
\hline
\textbf{Criterion} & \textbf{HD} & \textbf{D} & \textbf{C} & \textbf{P} & \textbf{F} \\
\hline

\textbf{Preparation \& Contribution} &
Insightful, advances discussion &
Regular, clear reasoning &
Occasional, appropriate &
Basic, prompted only &
Little or none \\
\hline

\textbf{Engagement with Others} &
Synthesises ideas, elevates group &
Builds on others’ ideas &
Responds to peers &
Listens, limited interaction &
Disengaged \\
\hline

\textbf{Curiosity \& Thinking} &
Explores ambiguity, asks deep questions &
Engages with complexity &
Asks genuine questions &
Attempts answers when prompted &
Avoids engagement \\
\hline

\textbf{Booklet Evidence} &
Complete, thoughtful, annotated &
Mostly complete, good detail &
Adequate, some gaps &
Minimal effort evident &
Sparse or missing \\
\hline

\end{tabularx}
}

\vspace{1em}

\subsection*{Overall Performance}
\vspace{0.5em}

\noindent
\fbox{HD} \hspace{1em}
\fbox{D} \hspace{1em}
\fbox{C} \hspace{1em}
\fbox{P} \hspace{1em}
\fbox{F}

\vspace{1.5em}

\subsection*{Notes (optional)}
\vspace{3em}
\noindent\rule{\textwidth}{0.4pt}

\vspace{1em}
\noindent\rule{\textwidth}{0.4pt}

\vspace{1em}
\noindent\rule{\textwidth}{0.4pt}